% Partie : sets-functions-relations
% Chapitre : sets
% Section : russells-paradox

\documentclass[../../../include/open-logic-section]{subfiles}


\begin{document}

\olfileid[fr]{sfr}{set}{rus}
\olsection{Le paradoxe de Russell}

L'extensionnalité justifie la notation $\Setabs{x}{\phi(x)}$ pour
\emph{l'}ensemble des $x$ tels que~$\phi(x)$. Mais ce qu'elle
justifie \emph{exactement}, c'est ceci : \emph{s'il existe}
un ensemble dont les membres sont tous les objets vérifiant
$\phi$, et eux seuls, \emph{alors} il n'y en a qu'un.
Autrement dit, une fois~$\phi$ fixé, l'ensemble
$\Setabs{x}{\phi(x)}$ est unique, \emph{s'il existe}.

Cette condition est essentielle !{} Toutes les propriétés ne
permettent pas de définir un ensemble par \emph{compréhension} :
certaines ne définissent \emph{aucun} ensemble.
Si toute propriété définissait un ensemble, nous aboutirions
à des contradictions. Le paradoxe de Russell en est l'exemple
le plus célèbre.

Les ensembles peuvent être des !!{element}s d'autres ensembles :
par exemple, l'ensemble des parties d'un ensemble~$A$ est
constitué d'ensembles. Il est donc légitime de se demander
si un ensemble est !!a{element} d'un autre.
Un ensemble peut-il être membre de lui-même ?
Rien dans l'idée intuitive d'ensemble ne semble l'exclure.
Si \emph{tous} les ensembles forment une collection d'objets,
on pourrait penser les rassembler en un seul ensemble :
l'ensemble de tous les ensembles. Étant lui-même un ensemble,
il serait !!a{element} de l'ensemble de tous les ensembles.

Le paradoxe de Russell apparaît lorsqu'on considère la propriété
de ne pas avoir soi-même pour !!{element}, autrement dit
de \emph{ne pas s'appartenir à soi-même}.
Supposons qu'il existe un ensemble de tous les ensembles
qui ne s'appartiennent pas à eux-mêmes. L'ensemble
\[
R = \Setabs{x}{x \notin x}
\]
existe-t-il ? On peut démontrer que non.

\begin{thm}[Paradoxe de Russell]\ollabel{thm:russells-paradox}
	Il n'existe pas d'ensemble $R = \Setabs{x}{x \notin x}$.
\end{thm}

\begin{proof}
Si $R = \Setabs{x}{x \notin x}$ existe, alors
$R \in R$ si et seulement si $R \notin R$,
ce qui est contradictoire.
\end{proof}

\begin{tagblock}{novice}
\begin{explain}
Reprenons cette démonstration plus en détail.
Si $R$ existe, on peut se demander si $R \in R$.
Supposons que $R \in R$. Or $R$ est défini comme
l'ensemble de tous les ensembles qui ne sont pas des
!!{element}s d'eux-mêmes. Si $R \in R$, alors $R$ ne vérifie
pas la propriété qui définit $R$. Mais seuls les ensembles
qui vérifient cette propriété appartiennent à~$R$.
Donc $R$ ne peut pas être !!a{element} de~$R$,
c'est-à-dire $R \notin R$. Mais $R$ ne peut pas à la fois
être et ne pas être !!a{element} de~$R$ :
nous obtenons une contradiction.

Puisque l'hypothèse $R \in R$ conduit à une contradiction,
on a $R \notin R$. Mais cela conduit aussi à une contradiction !{}
En effet, si $R \notin R$, alors $R$ vérifie lui-même
la propriété qui définit $R$. Ainsi, $R$ serait !!a{element} de $R$,
comme tous les autres ensembles qui ne s'appartiennent pas
à eux-mêmes. Là encore, il ne peut pas à la fois
ne pas être et être !!a{element} de~$R$.
\end{explain}
\end{tagblock}

\begin{digress}
Comment construire une théorie des ensembles qui évite
le paradoxe de Russell, c'est-à-dire qui n'affirme pas,
de façon \emph{contradictoire}, l'existence de
$R = \Setabs{x}{x \notin x}$ ?
Il faut poser des axiomes qui précisent exactement
dans quelles conditions les ensembles existent,
et dans quelles conditions ils n'existent pas.

La théorie des ensembles esquissée dans ce chapitre
ne le fait pas : elle est \emph{naïve}.
Elle dit seulement que les ensembles satisfont le principe
d'extensionnalité et qu'à partir d'ensembles donnés, on peut
former leur réunion, leur intersection, etc.
Il est possible de développer la théorie des ensembles
avec davantage de rigueur. \oliflabeldef{cumul:::part}{Ce sera
l'objet de la partie \olref[cumul][][]{part}.
Pour l'instant, nous poursuivrons de manière naïve,
en veillant à éviter les contradictions.}{}
\end{digress}


\end{document}
